\title{Parameter-Efficient Orthogonal Finetuning via Factorization}
For the special case of BOFT($1$,$d$), the model degenerates to the original OFT~\cite{qiu2023controlling} where the orthogonal matrix is fully dense and unconstrained. As our representation relies on expressing the orthogonal matrix as a sequence of sparse matrix multiplications, these multiplications must be executed at every training step. In what follows, we describe how adopting the butterfly structure furthers parameter-efficient representations. In contrast to the block-diagonal scheme employed by OFT to balance expressivity and regularization, BOFT leverages the butterfly structure to achieve a smoother transition between matrices sampled from the entire orthogonal group (offering expressivity) and the identity matrix (ensuring regularity). More formally, for a pretrained linear transformation $\bm{W}^0\in\mathbb{R}^{d\times n}$, OFT learns an orthogonal matrix $\bm{R}\in\mathbb{R}^{d\times d}$, altering the computation from $\bm{z}=(\bm{W}^0)^\top\bm{x}$ to $\bm{z}=(\bm{R}\bm{W}^0)^\top\bm{x}$, with $\bm{x}\in\mathbb{R}^{d}$ and $\bm{z}\in\mathbb{R}^{n}$ representing the input and output, respectively. To further economize parameters within our information transmission paradigm, we draw on the butterfly network architecture used in the Cooley-Tukey fast Fourier transform algorithm~\cite{cooley1965algorithm}, which facilitates rapid data exchanges among nodes~\cite{klasing1994broadcasting}. Our proposed information transmission viewpoint paves the way for integrating insights from the domain of computer networking, where network topology efficiency for information dissemination is extensively studied. Compared to approaches such as LoRA~\cite{hulora2022}, which rely on low-rank matrix structures, our method underscores that both low-rank and sparse matrices are foundational classes of structured matrices~\cite{candes2011robust} utilized for parameter-efficient modeling. Consequently, expressing $\bm{R}$ as a product $\bm{B}_m\bm{B}_{m-1}\cdots\bm{B}_1$ corresponds to staged information exchanges structured by the graph topology defined by each $\bm{B}_i$. The butterfly architecture, in conjunction with permutations, can exactly reproduce numerous well-known fast linear transforms~\cite{dao2019learning,dao2019kaleidoscope} (\emph{e.g.}, fast Fourier transform, Hadamard transform). Building on the generalization proposed in \cite{chen2022pixelated}, we define a block butterfly component $\tilde{\bm{B}}^b(d,k)$ in which each scalar entry in $\bm{d}_i$ (for any $i$) is replaced by a $b\times b$ matrix. This escalation has made training new foundation models from the beginning increasingly unfeasible for most research groups. While LoRA modifies the weights by adding a learnable low-rank component, OFT leverages a multiplicative orthogonal matrix for weight modification. To mitigate this, we analyze OFT through the lens of information flow, and establish several criteria that facilitate greater parameter efficiency. Standard Dropout~\cite{srivastava2014dropout} is naturally compatible with LoRA but does not directly apply to BOFT due to its multiplicative parameter updates. This analytical framework offers a means to assess OFT's parameter efficiency from a graph-based viewpoint, and it allows for the easy construction of feasible matrix factorizations. Notably, the sparsity structure in each butterfly component can be regarded as a permutation of the nonzero pattern found in $\tilde{\bm{B}}(d,2)$, thereby making the parameterization of every butterfly component as an orthogonal matrix straightforward. In finetuning approaches, models are generally initialized to exactly match the pretrained weights, in order to minimize deviation from the original performance. This approach can be considered as learning the bilinear similarity matrix $\bm{R}$ under a stringent constraint—namely, the orthogonality of $\bm{R}$. Figure~\ref{fig:info_trans} illustrates an instance of practical matrix factorization, which requires only $10$ edges overall—fewer than the single dense orthogonal case. While OFT has been shown to deliver notable improvements in generalization and convergence—especially in text-to-image diffusion model finetuning~\cite{qiu2023controlling}—it introduces a substantial number of learnable parameters owing to the dimensionality of the orthogonal transformations. However, LoRA applies a fixed rank across all layers, motivating later developments \cite{zhang2022adaptive,valipour2022dylora,zhang2023increlora} that dynamically assign rank per layer for more balanced parameter allocation. Employing butterfly factorization, one can synthesize a $d$-dimensional dense matrix through the product of $\mathcal{O}(\log d)$ sparse matrices, each possessing only $\mathcal{O}(d)$ nonzero elements. For instance, BOFT($9$,$2$) demonstrates superior expressivity to BOFT($1$,$16$) with substantially fewer parameters. OFT leads to better generalization and greater training stability relative to LoRA, especially in text-to-image diffusion model finetuning, yet typically entails more trainable parameters. A variety of PEFT strategies have been developed~\cite{houlsby2019parameter,aghajanyan2020intrinsic,hulora2022,edalati2022krona,wang2022adamix,gheini2021cross,zaken2022bitfit,guo2020parameter,sung2021training,ansell2022composable,lester2021power,li2021prefix,vu2022spot,he2021towards,mao2021unipelt,karimi2021compacter,liu2022few,sung2022lst,chen2022parameter,jia2022visual,chen2022adaptformer,zhang2022neural,jie2023fact,lian2022scaling,luo2023towards}, with reparameterization-based approaches~\cite{aghajanyan2020intrinsic,hulora2022,edalati2022krona} being especially pertinent to our research. However, advances in performance have been closely tied to steep growth in parameter counts—GPT-3, for example, incorporates approximately 175B parameters. 
\begin{abstract}
With the widespread adoption of large foundation models, the cost of training such models from the ground up has become increasingly prohibitive. Each block $\bm{B}^F_i(k)$, for all $i$, refers to a butterfly factor as given in Equation~\ref{eq:but_factor}. This property aligns with our setting, where information at each node in the input layer should propagate to all output nodes. In general, BOFT($m$,$b$) involves $\frac{1}{2}(b-1)dm$ effective trainable parameters when adapting a linear layer of shape $d\times n$. LoRA~\cite{hulora2022} implements finetuning by updating the pretrained weight matrices through the addition of the product of two low-rank matrices, showing strong results in natural language processing tasks. Such $d$ edges are a direct consequence of orthogonality and should be omitted from the count of trainable edges, since these entries are not adjustable (\emph{e.g.}, in a $d \times d$ orthogonal matrix with $d$ non-zeros, only values $\pm 1$ are possible). In general, BOFT achieves higher parameter efficiency than OFT. Such a block-diagonal sparsity configuration, although it economizes parameters, can bring about adverse inductive biases (\emph{e.g.}, diminished representational capacity and an inability to closely mimic standard linear mappings). Concretely, when synthesizing a dense matrix $\bm{R}\in\mathbb{R}^{d\times d}$ via a sequence of $m$ square matrices as $\thickmuskip=2mu \medmuskip=2mu \bm{R}=\bm{B}_m\bm{B}_{m-1}\cdots\bm{B}_1$, we interpret this process as the propagation of information across a $d\times (m+1)$ node grid, illustrated in Figure~\ref{fig:info_trans}. In this framework, all neurons within a layer are mapped via a shared orthogonal matrix, guaranteeing the invariance of pairwise angular relationships among neurons during finetuning. Bringing together these elements, the forward computation in BOFT is

\begin{equation}
    \begin{aligned}\nonumber
    \bm{z}=\big(\bm{R}(m,b)\cdot\bm{W}^0\big)^\top\bm{x},~~~~\text{s.t.}~\bigg{\{}\bm{R}(m,b)=\prod_{i=1}^m \tilde{\bm{B}}^b_{(d,i)}~~\&~~\underbrace{\big(\tilde{\bm{B}}^b_{(d,j)}\big)^\top\tilde{\bm{B}}^b_{(d,j)}=\tilde{\bm{B}}^b_{(d,j)}\big(\tilde{\bm{B}}^b_{(d,j)}\big)^\top\!=\bm{I}_d}_{\forall j\in [1, m]}\bigg{\}},
    \end{aligned}
\end{equation}

In the formulation above, the notation $\tilde{\bm{B}}^b(d,2^{m - i+1})$ is abbreviated as $\tilde{\bm{B}}^b_{(d,i)}$ for notational simplicity, and $\bm{I}_d$ denotes the $d \times d$ identity matrix. Beyond signal processing, the butterfly topology is also prevalent in computer network architecture~\cite{solihin2015fundamentals,li2011linear} as an effective means for information dissemination. For any entry $R_{ij}$ in $\bm{R}$, a zero signifies an absence of information transfer from the $j$-th node to the $i$-th, whereas a nonzero entry enables such transmission. Both OFT and BOFT implement orthogonal transformation by left-multiplying with $\bm{R}$, producing updated weights $\bm{R}\bm{U}\bm{\Sigma}\bm{V}^\top$. \end{itemize}

\section{Related Work}

\textbf{Parameter-efficient finetuning (PEFT)}. For ease of reference, BOFT using $\bm{R}(m,\frac{b}{2})$ is referred to as BOFT($m$,$b$), with $b\geq 2$. For improved parameter-efficiency, the block-diagonal pattern is achieved by structuring $\bm{R}$ as $\text{diag}(\bm{R}_1,\bm{R}_2,\cdots,\bm{R}_r)$, where each $\bm{R}_i\in\mathcal{R}^{b\times b}$ is an orthogonal matrix and $br=d$. The butterfly graph naturally lends itself to a sparse matrix factorization methodology known as butterfly factorization. However, the block-diagonal configuration, while economizing parameters, imposes a constraint: it implicitly assumes that the coordinates of a neuron's vector (i.e., the columns of $\bm{W}^0$) can be partitioned into $r$ groups, each transformed by separate orthogonal matrices. Our primary contributions are summarized below:

\begin{itemize}[leftmargin=*,nosep]
    \item We address the challenge of parameter efficiency in orthogonal finetuning via a new information transmission framework, which conceptualizes the construction of parameter-efficient dense orthogonal matrices as an information flow problem over a grid-based graph. These works generally employ the butterfly structure to reparameterize weight matrices directly, initializing and training neural networks from the ground up. Recognizing the inherently structured composition of weight matrices, a more systematic strategy focuses on conserving the relational structure among the weights—specifically, preserving the pairwise angles between neurons. Orthogonality, by its nature, demands fewer parameters than dense matrices and thereby introduces further dependencies among edges. Information is relayed sequentially, beginning with $\bm{B}_1$ and concluding at $\bm{B}_n$. & \bm{0} \\
    \bm{0}\!\!\!\!\! To do so, it suffices to ensure that each multiplicative component $\tilde{\bm{B}}(d,k)$ comprising the butterfly matrix $\bm{B}(d)$ remains orthogonal. \textbf{Butterfly structures}. Our central hypothesis is grounded in the observation that, whereas the butterfly structure is flexible enough to represent numerous well-known linear transformations, the block-diagonal format utilized in OFT cannot achieve this capability. \section{An Information Transmission View on Orthogonal Finetuning}

We begin by introducing essential concepts underlying OFT. To ensure orthogonality in the block butterfly component $\tilde{\bm{B}}^b(d,2)$, each $2b\times 2b$ matrix block is specifically parametrized to be orthogonal. We introduce an innovative strategy that composes the dense orthogonal matrix from several sparse, factorized matrices. Building upon connections between hyperspherical energy and generalization~\cite{liu2018learning,liu2021orthogonal}, \cite{qiu2023controlling} introduce orthogonal finetuning (OFT), which optimizes an orthogonal matrix to rotate the neurons within a layer. This prompts an important question: \emph{Is it possible to design a dense orthogonal matrix structure that retains parameter-efficiency?}

In response to this problem, we introduce an approach that constructs a dense orthogonal matrix by multiplying several sparse orthogonal matrices. To facilitate the discovery of suitable matrix decompositions, we develop a graphical information transmission perspective. Model finetuning is fundamentally guided by the aim of maintaining a resemblance between the finetuned model and its pretrained counterpart, as quantified by specific metrics, in order to retain the knowledge acquired during pretraining. Of these approaches, model finetuning stands out due to its simplicity, effectiveness, and the fact that it does not introduce additional inference latency. For instance, with a $2 \times 2$ orthogonal matrix, setting the $(1,1)$ element to zero forces the $(2,2)$ entry to be zero as well (\emph{i.e.}, the absence of one edge dictates the absence of another). Setting $d=2^N$, we recursively construct the $d$-dimensional \emph{butterfly matrix} $\bm{B}(d)\in\mathbb{R}^{d\times d}$ as

\begin{equation}
\begin{aligned}
    \!\!\bm{B}(d)=\tilde{\bm{B}}(d,d)\!\cdot\!\begin{bmatrix}
    \bm{B}_1(\frac{d}{2})\!\!\!\!\! By maintaining the orthogonality of every sparse matrix in this product, we establish an orthogonal parameterization with just $\mathcal{O}(d\log d)$ parameters, representing a substantial improvement over traditional dense formulations. Each colored curve tracks BOFT configurations sharing the same block size, with the leftmost point reflecting the error for BOFT($1$,$b$) (\emph{i.e.}, the basic block-diagonal OFT with block size $b$). If employing the butterfly architecture, specifically with $m=\log d$ and $b=2$, BOFT utilizes $\mathcal{O}(d\log d)$ parameters. To adapt a pretrained weight matrix, OFT expresses the updated weight matrix as a product of a trainable orthogonal matrix and the fixed, pretrained weight matrix. Beyond that, butterfly structures support fast matrix-vector multiplication~\cite{michielssen1996multilevel,o2010algorithm}, compact matrix representation~\cite{li2015butterfly}, and high-speed network transmission~\cite{klasing1994broadcasting,li2003linear,rajasingh2016transmission}. Additionally, it is still uncertain whether the butterfly network structure represents the optimal approach for information transmission. In contrast, standard OFT imposes a block-diagonal form on the orthogonal matrix~\cite{qiu2023controlling}, where smaller block sizes yield greater parameter savings. In this representation, each matrix $\bm{B}_i$ delineates connections from nodes at the $i$-th layer to those at the $(i+1)$-th layer; specifically, the $(j_1, j_2)$ element in $\bm{B}_i$ determines if there is a directed edge from node $j_2$ in the $i$-th layer to node $j_1$ in the subsequent layer (a zero here denotes absence of an edge). The specific inductive bias imparted by butterfly factorization is posited to enhance generalization performance, as its inherent structural patterns align with those found in numerous well-established linear transformations~\cite{dao2019kaleidoscope}, including the discrete Fourier transform, discrete sine/cosine transform, and Hadamard transform. \textbf{Comparison to butterfly-based sparse training}. Strategies for reducing the training overhead of BOFT remain to be explored. Figure~\ref{fig:info_trans} illustrates this for the case $\bm{R}=\bm{B}_5\bm{B}_4\bm{B}_3\bm{B}_2\bm{B}_1$, with both the sparsity structures and the resulting graph shown. \item We extend the application of orthogonal finetuning, for the first time, to a variety of domains beyond the original setting of controllable text-to-image generation~\cite{qiu2023controlling}, demonstrating its promise as a versatile approach to model finetuning. The results further affirm BOFT's advantages as a general-purpose finetuning technique. In BOFT, $\bm{R}(m,b)$ is commonly chosen from a structured subgroup of the orthogonal group, restricting the overall hypothesis space and thereby introducing inductive bias. Consequently, it is vital to develop effective techniques for customizing these models for new tasks. \item We offer theoretical justifications for BOFT’s effectiveness in drastically reducing trainable parameters without sacrificing expressivity or stability during optimization. LoRA achieves parameter efficiency by constraining its updates to a low-rank structure. \end{document} Importantly, following the finetuning process, the orthogonal matrices optimized by BOFT can be merged directly into the pretrained model, ensuring that inference incurs no additional latency. Experimental evaluations substantiate the effectiveness of BOFT across various domains, including large language models, large-scale vision models, and text-to-image generative models. \begin{theorem}[Expressivity of BOFT]\label{thm:exp_boft}
    BOFT possesses greater expressive capacity than OFT when matched for block size. Nevertheless, BOFT is not without its drawbacks. Our objective, therefore, is to synthesize a dense orthogonal matrix from $m$ sparse orthogonal factors, all while minimizing the number of trainable parameters employed. For butterfly components $\tilde{\bm{B}}^b(d,k)$ with $k>2$, the nonzero patterns represent block-wise permutations of that in $\tilde{\bm{B}}^b(d,2)$, thus facilitating the extension to orthogonal matrices in a parallel manner. \section{Intriguing Insights and Discussions}

\textbf{Expressivity of BOFT}. Furthermore, the optimal selection of such sparsity patterns remains an unresolved challenge. \section{Orthogonal Parameterization by Butterfly Factorization}

The butterfly structure, foundational to the Cooley-Tukey fast Fourier transform algorithm, enables the efficient computation of the Fourier transform, wherein a small perturbation in the frequency domain leads to broad changes in the spatial domain. To address this, we introduce a multiplicative Dropout scheme specifically designed for BOFT. Let $k \geq 2$ be a power of $2$; we define the \emph{butterfly factor} $\bm{B}^F(k)$ as follows:

\begin{equation}\label{eq:but_factor}
\begin{aligned}
    \bm{B}^F(k)=\begin{bmatrix}
    \text{diag}(\bm{d}_1) & \text{diag}(\bm{d}_2) \\
    \text{diag}(\bm{d}_3) & \text{diag}(\bm{d}_4)
    \end{bmatrix}\in\mathbb{R}^{k\times k},
\end{aligned}
\end{equation}

where each $\bm{d}_i\in\mathbb{R}^{\frac{k}{2}}$ for $i=1,\ldots,4$ are parameter vectors. The orthogonality of $\tilde{\bm{B}}(d,2)$ can be readily imposed via the Cayley transform, or equivalently, by using $2$-dimensional rotation matrices to parametrize each $2\times 2$ block, consistent with approaches such as in \cite{liu2021orthogonal,qiu2023controlling}. At this stage, we can leverage the butterfly matrix to represent any orthogonal matrix by appropriately selecting parameters. LoRA~\cite{hulora2022} incorporates Dropout layers within its low-rank updates to combat overfitting. Leveraging this structure, we present Orthogonal Butterfly~(BOFT), a new, highly parameter-efficient finetuning approach. Owing to the ease of use, such low-rank weight reparameterizations are widely adopted~\cite{dettmers2023qlora,zi2023delta,chavan2023one}. This distinction is visually summarized in Figure~\ref{fig:comp_lora}. With BOFT($\log \frac{2d}{b}$,$b$), the block butterfly matrix $\bm{B}^{\frac{b}{2}}(d)$ functions as $\bm{R}$, yielding a full orthogonal matrix with dense structure. Comparatively, the original OFT with a dense orthogonal matrix requires $\mathcal{O}(d^2)$ parameters, while the block-diagonal OFT with $r$ blocks has a parameter complexity of $\mathcal{O}(bd)$. \textbf{BOFT Multiplicative Dropout}. Prior studies~\cite{dao2019learning,dao2019kaleidoscope,chen2022pixelated,dao2022monarch} have investigated sparse training methods utilizing butterfly parameterizations. Overall, the butterfly matrix models a more constrained subset within the orthogonal group than the block-diagonal approach, and thus, BOFT is provably more expressive than OFT at a fixed block size. This work explores a systematic approach to task adaptation called Orthogonal Finetuning (OFT). In Figure~\ref{fig:blk_diag}, the block-diagonal configuration of $\bm{R}$ in the standard OFT setup is depicted. For example, the configuration shown in Figure~\ref{fig:info_trans} reaches dense connectivity with $10$ edges, while a butterfly network requires only $8$. Similar correspondences between graph reachability and matrix sparsity hold for further truncated products such as $\bm{R}=\bm{B}_3\bm{B}_2\bm{B}_1$ and $\bm{R}=\bm{B}_2\bm{B}_1$. Nevertheless, it is not yet fully understood how closely our orthogonal butterfly matrix can approximate an arbitrary orthogonal matrix. BOFT expands the tunable parameter space by embedding structural priors, thereby enabling the identification of an optimal compromise between expressiveness and regularization. A notable feature shared by both block-diagonal and butterfly constructions is their inherent sparsity, a property that enables both schemes to reduce the number of required trainable parameters in orthogonal matrices. To tackle this issue, our objective is to construct a dense orthogonal matrix without incurring high parameter costs. Striking an effective balance between model expressivity and inductive regularity is crucial to ensuring strong generalization in model finetuning. It is important to highlight that orthogonality imposes nuanced constraints on the edges linking successive levels. This analogy enables the creation of diverse sparse matrix factorization schemes that curtail trainable parameters while retaining sufficient expressiveness for dense matrix construction. We systematically assess the performance of BOFT across various settings, adapting large-scale vision transformers, language models, and text-to-image diffusion models to an array of downstream vision and language tasks. Efficiently repurposing these models for downstream applications is thus essential. This, in turn, is likely to enhance generalization and reduce the risk of overfitting. To showcase this, BOFT is tested across several adaptation scenarios, spanning computer vision and natural language processing. To approximate any orthogonal matrix of dimension $d$, one can construct a product of butterfly matrices:  $\bm{B}_{d-1,1}(d)\bm{B}^\top_{d-1,2}(d)\cdots\bm{B}_{1,1}(d)\bm{B}^\top_{1,2}(d)$, with each $\bm{B}_{i,j}(d),\forall i,\forall j$ representing a butterfly matrix. The rapid expansion in size and capability of foundation models (\emph{e.g.}, \cite{brown2020language,radford2021learning,rombach2022high,kirillov2023segment}) has spurred research into methods that allow efficient finetuning of such models for downstream applications~\cite{treviso2023efficient,ding2023parameter,lialin2023scaling}. Restricting to $\bm{R}=\bm{B}_4\bm{B}_3\bm{B}_2\bm{B}_1$, i.e., connections from layer one to layer five, reveals that, for instance, node $1$ cannot reach node $3$, and so the entry at $(3,1)$ in the resulting matrix is zero. \input{rebuttal-results}

\section{Concluding Remarks and Limitations}

In this work, we introduce Orthogonal Butterfly, a versatile and parameter-efficient method for finetuning foundation models that leverages the butterfly structure. While OFT exhibits strong generalization capabilities, it requires updating a significant number of parameters due to the inherent dimensionality of orthogonal matrices. By integrating this parameterization into OFT, we present a new finetuning method named Orthogonal Butterfly (BOFT), which generalizes the standard OFT framework by encompassing it as a specific instance. Our empirical results indicate that BOFT significantly exceeds current leading methods, underscoring both its parameter efficiency and capacity for generalization. For example, while full fine-tuning is maximally expressive, it frequently produces less desirable performance. \textbf{Orthogonal finetuning as learning bilinear similarity}. Conventionally, forming a dense $d$-dimensional orthogonal matrix demands $\mathcal{O}(d^2)$ parameters, rendering parameter-efficiency elusive. Additional mathematical details are presented in Appendix~\ref{app:sn_property}. \textbf{Spectral properties}. Zero initialization is ensured by setting $\bm{R}$ to the identity matrix at the outset. BOFT encompasses OFT as a specific instance, offering a comprehensive orthogonal finetuning framework. There exist three principal strategies for such adaptation: \emph{model finetuning}~\cite{hulora2022,zaken2022bitfit,guo2020parameter,raffel2020exploring,qiu2023controlling,zhang2022adaptive,chavan2023one}, which tunes a subset of the model's original parameters without altering its architecture; \emph{adapter tuning}~\cite{rebuffi2017learning,houlsby2019parameter,pfeiffer2020adapterfusion,lin2020exploring,he2021towards}, which inserts additional, trainable modules into the network; and \emph{prompt tuning}~\cite{li2021prefix,lester2021power}, which adds learnable prefix tokens to model inputs. Similarly, BOFT($6$,$4$) attains comparable approximation errors to BOFT($2$,$16$) while requiring significantly fewer parameters. The outcomes presented in Figure~\ref{fig:para_eff} are averaged over 10 different random seeds. The perspective of information transmission provides valuable insights into the sparsity characteristics of factorization matrices within OFT. Enhancing the parameter-efficiency of OFT, therefore, becomes an important objective, particularly as its adaptability to broader task types remains unexplored. Orthogonal finetuning typically achieves superior spectral characteristics compared to LoRA, primarily due to its exact preservation of the spectral norm from the pretrained weight matrix $\bm{W}^0$. Additionally, the factorized matrix representation used in BOFT facilitates an interpretable interpolation of weights (refer to Figure~\ref{fig:boft_interpolation}). Despite empirical gains, this grouping—merely based on index positions—lacks conceptual justification, indicating that dense orthogonal matrices may be preferable. \end{theorem}

Based on Theorem~\ref{thm:exp_boft}, a straightforward generalization of BOFT emerges: the final orthogonal transformation is extended as $\thickmuskip=2mu \medmuskip=2mu \bm{R}^G(m_1,b_1,m_2,b_2,l)=\bm{R}_{l,1}(m_1,b_1)\bm{R}_{l,2}^T(m_2,b_2)\cdots\bm{R}_{1,1}(m_1,b_1)\bm{R}_{1,2}^T(m_2,b_2)$, with $\bm{R}_{i,j}^T(m,b)$ representing the orthogonal matrix type employed in BOFT. Within this context, producing a dense orthogonal matrix by the successive multiplication of sparse factors is analogous to broadcasting information across a grid-based network. The radix-2 Cooley-Tukey algorithm recursively breaks an $N$-point discrete Fourier transform into a pair of $\frac{N}{2}$-point transforms, resulting in a structure that can be expressed as a product of several sparse matrices (collectively referred to as a butterfly matrix). Therefore, to achieve a dense orthogonal matrix, the standard OFT necessitates choosing $b=d$ as the block size, whereas BOFT can generate a dense structure for any $b$. In contrast, \cite{dao2022monarch} addresses finetuning by initially projecting pretrained weights onto variants of butterfly matrices before updating the resulting components for specific downstream applications. This perspective is motivated by the idea that a dense $d$-dimensional square matrix can be seen as specifying dense pathways from one set of $d$ nodes to another. Butterfly matrices have been explored for parameterizing orthogonal matrices to streamline Gaussian elimination~\cite{parker1995random}, for improving the stability of recurrent neural network training~\cite{jing2017tunable}, and for efficient kernel approximations~\cite{munkhoeva2018quadrature}. It is important to highlight, however, that enhanced expressivity does not necessarily translate into improved fine-tuning outcomes. Drawing inspiration from the Cooley-Tukey fast Fourier transform, which excels at compactly transmitting information, we introduce an efficient orthogonal parameterization using butterfly structures. The depicted graph arises from unfolding the multiplication of these matrices. If we select $m_1=m_2=\log d$, $b_1=b_2=2$, and $l=d-1$, then $\bm{R}^G(m_1,b_1,m_2,b_2,l)$ spans the entire orthogonal group, aligning with what is known as the kaleidoscope hierarchy~\cite{dao2019kaleidoscope}. Since BOFT’s final orthogonal matrix arises from the product of several orthogonal matrices, training incurs a slightly higher computational cost relative to OFT. Given that $\bm{R}(m,b)$ consists of $m$ butterfly matrices, each easily expressible as $2b\times 2b$ block-diagonal orthogonal matrices, this Dropout randomly selects $p_1$ percent of the butterfly components and $p_2$ percent of the diagonal blocks within each component, then replaces the chosen blocks with identity matrices. This formulation has inherent links to concepts such as distance metric learning~\cite{xing2002distance} and the bilinear form~\cite{roman2005advanced}. \textbf{Inductive bias and generalization in BOFT}. If $\thickmuskip=2mu \medmuskip=2mu m=1$, BOFT($1$,$b$) becomes the block-diagonal OFT~\cite{qiu2023controlling} where the block size is $b$. Accordingly, any specific pattern of edges may be altered—either by necessitating or precluding additional edges—due to orthogonality. Specifically, to guarantee that each $\bm{B}_i$ is full-rank—a prerequisite for orthogonality—one must establish $d$ edges to set up a bijective mapping between nodes at the $i$-th and $(i=1)$-th levels; this ensures that the inter-level edge count is at minimum $d$ (\emph{e.g.}, see the $4$ edges in Figure~\ref{fig:info_trans}). By examining the structure of non-zero entries in the resulting orthogonal matrix, the information transmission viewpoint proves especially beneficial in the context of OFT, since only the configuration of non-zero, trainable elements in $\bm{R}$ is relevant, not their individual values. BOFT advances this goal by employing butterfly factorization to improve OFT’s parameter efficiency, thereby extending the applicability of orthogonal finetuning to a wide array of adaptation challenges. This approach allows us to constrain our search space to a smaller subset within the orthogonal group tailored for downstream applications. The butterfly structure, commonly found in fast linear operators like the discrete Fourier and cosine transforms, motivates our claim that the structured design of BOFT inherently promotes parameter efficiency by instilling an implicit inductive bias. More generally, for a dense matrix $\bm{R}\in\mathbb{R}^{d\times d}$ to arise, the collection $\bm{B}_m,\ldots,\bm{B}_1$ must generate a directed edge structure on a $d\times(m+1)$ grid such that each directed edge links nodes in adjacent columns (levels), ensuring all starting nodes at the first level can propagate information to every destination node at the $(m+1)$-th level. Within this framework, we find that the butterfly architecture offers a compelling solution for constructing sparse orthogonal matrix factorizations. Additionally, \cite{dao2019learning,dao2019kaleidoscope} employ butterfly parameterizations to devise efficient linear transformations, while \cite{chen2022pixelated,dao2022monarch} leverage these structures for resource-efficient neural network training. The rationale for employing an orthogonal transformation during fine-tuning is to retain the angular relationships between neurons~\cite{liu2018learning,liu2021orthogonal,qiu2023controlling}, thereby largely maintaining the semantic information encoded during pretraining. Our central contribution to enhancing parameter efficiency lies in representing dense orthogonal matrices as products of multiple sparse orthogonal matrices. We draw on the well-established butterfly graphs~\cite{cooley1965algorithm} from the Cooley-Tukey algorithm, which facilitate efficient full connectivity between $d$ sources and $d$ targets with just $\mathcal{O}(d \log d)$ edges. BOFT, in contrast, initializes every butterfly component as an identity matrix (i.e., initializing the skew-symmetric parameter matrix as all zeros in the Cayley transform). &\bm{B}_2(\frac{d}{2})
    \end{bmatrix}= \tilde{\bm{B}}(d,d)\tilde{\bm{B}}(d,\frac{d}{2})\cdots\tilde{\bm{B}}(d,2),
\end{aligned}
\end{equation}

Here, $\bm{B}_1(\frac{d}{2})$ and $\bm{B}_2(\frac{d}{2})$ denote butterfly matrices of dimension $\frac{d}{2}$. We anticipate that network architectures superior to the butterfly structure may ultimately yield even more efficient constructions of dense orthogonal matrices. As a parallel, LoRA uses a zero start for its low-rank matrices. We introduce the notion of the \emph{butterfly component} as $\thickmuskip=2mu \medmuskip=2mu \tilde{\bm{B}}(d,k)=\text{diag}(\bm{B}^F_1(k),\ldots,\bm{B}^F_{d/k}(k))$, where this is a block-diagonal matrix with size $d\times d$ composed of blocks each of size $k$. \end{abstract}

\section{Introduction}

State-of-the-art models, including ChatGPT~\cite{brown2020language,chen2021evaluating} and Stable Diffusion~\cite{rombach2022high}, exemplify the exceptional generalization achieved by large foundation models. BOFT, however, introduces a fundamentally distinct approach to finetuning by transforming model weights through the application of layer-shared weight matrices. \item Drawing inspiration from the butterfly decomposition as used in the Cooley-Tukey algorithm, we introduce Orthogonal Butterfly—a novel orthogonal finetuning technique grounded in the aforementioned information transmission framework, specifically designed for improved parameter efficiency. To investigate this, we perform an experiment attempting to approximate a randomly sampled dense orthogonal matrix~\cite{anderson1987generation} of size $512\times 512$. A straightforward implementation of dense connections between two layers involves $\mathcal{O}(d^2)$ edges (i.e., one fully connected orthogonal matrix), resulting in $d^2 - d$ actual edges; for instance, when $d=4$, this totals 12 edges. To preserve the orthogonality constraint on $\bm{R}$ throughout fine-tuning, we adopt the Cayley parameterization as proposed in \cite{qiu2023controlling,liu2021orthogonal}: $\bm{R}=(\bm{I}+\bm{Q})(\bm{I}-\bm{Q})^{-1}$, where $\bm{Q}$ is a skew-symmetric matrix satisfying $\bm{Q}=-\bm{Q}^\top$. For example, commonly used finetuning strategies typically involve selecting a modest learning rate, a heuristic that restricts the extent of deviation from the pretrained parameters. To systematically and intuitively investigate how sparsity emerges in orthogonal matrix decompositions, we recast the challenge of generating a dense orthogonal matrix in OFT as analogous to an information transmission scenario. The y-axis corresponds to the approximation error, whereas the x-axis represents the effective number of trainable parameters. \textbf{BOFT Identity Initialization}. In addition, the use of sparse matrix factorization within BOFT may give rise to additional implicit inductive biases~\cite{gunasekar2017implicit,li2018algorithmic,liu2019neural}. Departing from prior applications, this work centers on utilizing butterfly structures to improve the parameter efficiency within OFT. This method is driven by the realization that reducing the parameter count is possible by accepting increased computational cost in exchange for lower memory usage. Maintaining this property has been demonstrated to improve both the stability of training and the generalization ability of the model~\cite{miyato2018spectral,yoshida2017spectral}. This operation leaves the maximum singular value, and thus the spectral norm of $\bm{W}^0$, unchanged. Nevertheless, this increased parameter efficiency entails certain drawbacks: the orthogonal matrix is restricted to a predetermined sparse structure, and each block undergoes transformation independently. The BOFT framework can alternatively be interpreted as a method for learning bilinear similarity measures, represented by $\bm{w}^0_i\bm{R}\bm{x}$, where $\bm{w}^0_i$ refers to the $i$-th neuron (specifically, a column vector) of $\bm{W}^0$. Consequently, this yields an efficient framework for representing orthogonal matrices through products of multiple $2\times 2$ orthogonal blocks. The orthogonal matrix $\bm{R}(m,b)\in\mathbb{R}^{d\times d}$ is constructed as a sequential product of several orthogonal butterfly matrices. To mitigate this, OFT employs a block-diagonal configuration for the orthogonal matrices, thereby reducing parameter count. Although this structure contributes to parameter reduction, it is incapable of yielding a dense $\bm{R}$ matrix. Building on this observation, Orthogonal Finetuning~(OFT)~\cite{qiu2023controlling} has been developed. To elucidate, consider the special case of $\tilde{\bm{B}}(d,2)$, which has block size $2$. This preservation can be clearly observed through singular value decomposition, $\bm{W}^0=\bm{U}\bm{\Sigma}\bm{V}^\top$, where $\bm{U}$ and $\bm{V}$ denote orthogonal matrices and $\bm{\Sigma}$ is a diagonal matrix containing singular values. Advancing the field now relies on methods to repurpose pre-existing models for new objectives, bypassing full retraining. Interpreted through the lens of information transmission, the core requirements of our objective are: \emph{(i)} \textbf{dense connectivity}, ensuring every node at the input level maintains at least one pathway to each node at the output level, and \emph{(ii)} \textbf{minimal free edges}, aiming for the lowest possible total edge count that upholds orthogonality. For the entire product $\bm{B}_5\bm{B}_4\bm{B}_3\bm{B}_2\bm{B}_1$ to produce a dense result, all nodes in the initial layer must have a path to every node in the sixth layer. To clarify the matrix factorization process, we recast the generation of dense orthogonal matrices as an information dissemination challenge.